% Pursuing Stacks (cotes 134-2 to 134-8): archive cote 134-7, site page 50 (PDF page 51), typescript page 492.
% typescript. Draft compared: ps6.tex:2458-2524.
% AI-transcribed by Claude Opus 5.5 (claude-opus-5-5, adaptive thinking, effort high), Message Batches API, batch ps3 (msgbatch_01XdFmuVVBAcgc7EFxdtCqsj), 2026-11-12; not yet checked by eye.
\scanpage{134-7/50}
\mspage{492}

10.9.

117. Maybe the best will be to write up (and possibly develop some)
my reflections \xout{after the notes}\hand{(}of the two days after I stopped with the
notes\hand{)} roughly in the order as they occured.

There were some somewhat technical afterthoughts. One was about the
logical difficulty coming from the site Aff$_{/k}$ on page 446 not being an
$\underline{U}$-site (where $\underline{U}$ is the universe we are working in), hence strictly spea-%
king the category of all sheaves on this site is possibly not even a
$\underline{U}$-category (i.e. the Hom objects need not be small i.e. with cardinal in
$\underline{U}$ ), still less a topos, and hence the standard panoply of notions and
constructions in a topos does not apply. This does'nt look really serious
though, one way out is to limit beforehand the ``size'' of the unipotent
bundles we are allowing i.e. of the k-modules describing them, in terms
say of cardinality of a family of generators of the latter - and then
restrict accordingly the size of the k-algebras k' taken as ``arguments''
for our sheaves i.e. as objects of the basic site we\inserted{are} work\over{in.}{ing}\inserted{on.} For instanc\hand{e,}
when working with unipotent bundles of finite type only (i.e. correspon-%
ding to k-modules of finite type - a rather interesting and natural
finiteness \xout{xxqxx} condition\inserted{anyhow} on the components of a schematic model $X_{\ast}$ ),
it is appropriate to work on the ``fppf site'',where the arguments k' are
k-algebras of \uline{finite presentation}. If we should be unwilling to be limi\hand{\unsure{ted}}
by a fixed size restriction on the unipotent bundles we are working with
(and hence also on the corresponding homotopy types), we may have to work
with a hierarchy of size restrictions and passage from one to any other
less stringent one - a technical nuisance to be sure, if we do'nt find
a more elegant way out, but surely not a substantial difficulty. At the
present heuristic stage of reflections, it does'nt seem worth while real-%
ly to dwell on such questions any longer.

Another afterthought is about the functor
\[ M \longmapsto W(M) \quad : \quad \mathrm{Ab}_k =(k\text{-Mod}) \longrightarrow (\underline{O}_k\text{-Mod}) \]
from k-modules to \xout{quasix} $\underline{O}_k$-Modules - a fully faithful functor we know,
whose essential image by definition consists of the so-called \uline{quasi}-\uline{cohe-}%
\uline{rent} sheaves. A little caution is needed, as this functor is right
exact, but \uline{not} \uline{exact} i.e. it does not commute with formation of kernels,
because for a k-algebra k' which is'nt flat, the functor
\[ M \longmapsto M\otimes_k k' \]

\margin{50}
